\subsection{R8: the recovered Fermat constructions and their exact scope}

We independently reconstruct the section constructions in the recovered
\texttt{LM6} manuscript, E.16 and E.31--E.34.  Throughout this subsection
\[
 X_n:\ X^n+Y^n+Z^n=0,\qquad
 U_n=X_n\setminus\{XYZ=0\},\qquad
 t=-X^n/Z^n.
\]
The sign and the ground field are part of the assertion.  In particular,
the even plus curve is not the usual minus Fermat curve over
\(\mathbb Q\).  In a quotient, ``cuspidal'' initially means the image
of an actual cuspidal packet.  E.99 proves that in the full metabelian
quotient this agrees with lying, up to geometric conjugacy, in the
arithmetic normalizer of boundary inertia.  This agreement must not be
assumed for an abelian quotient, where its geometric group centralizes
every boundary line.

\begin{result}{E.98}{Proved}{Actual pro-prime and full pronilpotent Fermat sections over \(\mathbb Q\)}
Let \(p\in\{3,5\}\), \(r\geq1\), and \(n=p^r\).  Put
\(O=[0:-1:1]\), take the rational unit tangent with local parameter
\(s=-X/Z\), and write
\[
 N=\pi_1(U_{n,\overline{\mathbb Q}}),\qquad
 P=N^{(p)},\qquad V=P^{\rm ab},\qquad S=\{p,\infty\}.
\]
Every class of \(H^1(G_{\mathbb Q,S},V)\) lifts to a section of the
full geometrically pro-\(p\) extension of \(U_n/\mathbb Q\).
There are infinitely many such classes whose images are noncuspidal
in every positive-class intrinsic pro-\(p\) nilpotent quotient and
every positive-length intrinsic pro-\(p\) solvable quotient.
Moreover, these classes give actual noncuspidal sections of the full
geometrically pronilpotent extension, simultaneously at all coefficient
primes and all nilpotency classes.  No full profinite section is asserted.
\end{result}
\begin{proof}
We first specify the arithmetic quotient.  Let \(K=\mathbb Q(\mu_p)\)
and let \(\Omega\) be its maximal pro-\(p\) extension unramified outside
the unique prime above \(p\).  This extension is Galois over
\(\mathbb Q\).  Set
\[
 P_{\rm ar}=\operatorname{Gal}(\Omega/K),\qquad
 \Gamma=\operatorname{Gal}(\Omega/\mathbb Q),\qquad
 1\longrightarrow P_{\rm ar}\longrightarrow\Gamma
 \longrightarrow\Delta\longrightarrow1,
 \quad |\Delta|=p-1.                                      \tag{E.98.1}
\]
Here \(\Omega\) is a maximal pro-\(p\) extension of \(K\), not the
maximal unrestricted extension unramified outside \(p\).

For clarity, the arithmetic freeness in this small-prime application
can be checked directly.  The Eisenstein integer ring is Euclidean,
and the Minkowski bound for \(\mathbb Q(\zeta_5)\) is
\[
 \frac{4!}{4^4}\left(\frac4\pi\right)^2\sqrt{125}<2.
\]
Thus both fields have class number one.  With
\(r_2=(p-1)/2\), the \(S_K\)-unit theorem and Kummer theory give
\(\dim_{\mathbb F_p}H^1(G_{K,S_K},\mathbb F_p)=r_2+1\).
The global Euler formula gives
\(\#H^0\#H^2/\#H^1=p^{-r_2}\), hence
\(H^2(G_{K,S_K},\mathbb F_p)=0\).
The precise finite-cohomology inputs are
\citet[I.4.15, p.~62; I.5.1, pp.~66--67]{CohMilneADT}.
If \(J\) is the kernel of \(G_{K,S_K}\to P_{\rm ar}\), a nonzero
invariant character \(J\to\mathbb F_p\) would give a strictly larger
pro-\(p\) quotient.  The five-term sequence therefore injects
\(H^2(P_{\rm ar},\mathbb F_p)\) into the zero group just computed.
The pro-\(p\) freeness criterion shows that \(P_{\rm ar}\) is free of
rank \(r_2+1\).  This is also the applicable special case of
\citet[Lemma 3.2, author p.~14]{LitHoshi2010}.
Restriction and averaging over \(\Delta\) give
\(\operatorname{cd}_p\Gamma\leq1\).

Next we check the actual action.  The cover of the tripod has, at the
chosen tangential point, the expansions
\[
 t=s^n,\qquad -Y/Z=(1-s^n)^{1/n},                         \tag{E.98.2}
\]
with leading constants one.  The induced Puiseux fibre functors are
therefore \(G_{\mathbb Q}\)-equivariant with the unit tangent at zero
on the tripod.  Consequently \(P\) is the actual open subgroup
\(\ker(F^{(p)}\to(\mathbb Z/n)^2)\), where \(F\) is its geometric
free group.  The pro-\(p\) topology on this subgroup is the full one:
intersecting the finitely many conjugates of a subgroup of \(p\)-power
index leaves a normal subgroup with \(p\)-group quotient.
Schreier's formula gives \(\operatorname{rank}P=n^2+1\).
The unit-tangent action on \(F^{(p)}\) is unramified away from \(p\)
by the pointed specialization theorem
\citep[Definition 2.0 and Theorem 2.1, pp.~198--200]{AuditWojtkowiak2007}.
Its residual action over \(K\) is trivial, and the kernel of
\(\operatorname{Aut}(F^{(p)})\to\operatorname{GL}_2(\mathbb F_p)\)
is pro-\(p\).  Thus this pointed action, and its restriction to
\(P\), factor through the actual quotient (E.98.1).

A continuous \(V\)-cocycle on \(G_{\mathbb Q,S}\) factors through
\(\Gamma\): over \(K\) its graph has pro-\(p\) image in
\(V\rtimes\rho(G_{K,S_K})\), so kills \(J\).
The pro-\(p\)-kernel lifting criterion applied to
\[
 P\rtimes\Gamma\longrightarrow V\rtimes\Gamma
                                                               \tag{E.98.3}
\]
lifts that graph.  We use precisely the general criterion in
\citet[Proposition 4.3, p.~149]{LitPachecoStevenson2000}, not the
positive-characteristic geometric theorems elsewhere in that paper.
The lifted homomorphism projects identically to \(\Gamma\), and its
cocycle inflates to \(G_{\mathbb Q}\).  This solves the lifting
problem over \(\mathbb Q\) itself; it does not descend an arbitrarily
chosen section over \(K\).

It remains to produce the noncuspidal classes.  The generalized
Jacobian gives the integral exact sequence
\[
 0\longrightarrow W\longrightarrow V\longrightarrow M
 \longrightarrow0,
 \qquad M=T_p\operatorname{Jac}(X_n),                    \tag{E.98.4}
\]
where \(W\) is the boundary torus module.  For every finite
\(p\)-primary module with action through \(\Gamma\), its restriction
to \(G_{K,S_K}\) has a filtration with trivial \(\mathbb F_p\)
factors.  The preceding \(H^2\)-vanishing and averaging show that its
\(H^2(G_{\mathbb Q,S},-)\) vanishes.  Finiteness of the corresponding
\(H^1\) groups removes the derived-limit term for each fixed lattice.
In particular \(H^2_{\rm cts}(G_{\mathbb Q,S},W)=0\), so (E.98.4)
surjects on \(H^1\).

For later use, the same argument yields the rank formula
\[
 \operatorname{rank}_{\mathbb Z_p}H^1(G_{\mathbb Q,S},L)
       =\dim_{\mathbb Q_p}(L\otimes\mathbb Q_p)^-        \tag{E.98.5}
\]
for a finite free lattice \(L\) with this action and no rational
invariants.  Indeed \(H^2_{\rm cts}=0\); applying the finite Euler
formula to \(L/p^mL\), then taking the coefficient of \(m\) in its
\(p\)-adic order, proves (E.98.5).  Finite generation follows as
well by restricting cocycles to the finitely generated free group
\(P_{\rm ar}\).  For \(M\), the invariant group is zero by
Mordell--Weil, and complex conjugation has \(g\) eigenvalues of each
sign because it reverses the intersection pairing.  Here
\(g=(n-1)(n-2)/2>0\).  Thus \(H^1(G_{\mathbb Q,S},M)\) has rank
\(g\).  Choose infinitely many distinct infinite-order classes,
lift them through (E.98.4), and then through (E.98.3).

The ratio of the tangent-line equations at any two Fermat cusps
\(d,O\) has divisor \(n(d-O)\): each such line meets \(X_n\) only
at that cusp, with multiplicity \(n\).  Consequently every cusp
class in the proper Jacobian is \(n\)-torsion.  An actual cuspidal
section has that proper Kummer class, independently of its tangent.
Our chosen proper class has infinite order.  This detects
noncuspidality already in the proper abelian quotient, hence in
every positive nilpotent or solvable quotient under consideration.
Distinct cusp lines remain distinct in geometric abelianization;
an arithmetic normalizer projecting onto \(G_{\mathbb Q}\) must
therefore lie over one of the three rational cusps.

Finally, the maximal pronilpotent quotient of any profinite group is
canonically the product of its maximal pro-prime quotients.  A finite
nilpotent quotient is the product of its Sylow subgroups; conversely,
the image in any finite collection of distinct-prime quotients is
surjective by the coprime subdirect-product argument.  Compactness
passes from those finite products to the whole product.  Applying
this also in each fixed nilpotency class gives
\[
 N^{\rm nil}=\prod_qN^{(q)},\qquad
 N/\gamma_{c+1}N=\prod_qN^{(q)}/\gamma_{c+1}N^{(q)}.
                                                               \tag{E.98.6}
\]
These identifications are equivariant for the actual action.  Place
one of the constructed cocycles in the \(p\)-factor and the identity
cocycle in every other factor.  This is one continuous product-valued
cocycle, not a collection of unrelated levelwise choices.  Its proper
\(p\)-abelian class still excludes every cusp.  This proves the
full pronilpotent assertion and reconstructs LM6 E.16 and E.32.
\end{proof}

\begin{result}{E.99}{Proved}{Full profinite metabelian zero-abelian Fermat families}
With \(p,n,N,O\) as in E.98, put
\[
 M_{\rm all}=N/N'',\qquad A=N^{\rm ab},\qquad
 d_n=\frac{(n^2-1)(n^2+3)}4.
\]
There is a continuous injective family
\[
 \mathbb Z_p^{d_n}\longrightarrow
 \operatorname{Sec}(M_{\rm all}\rtimes G_{\mathbb Q})/M_{\rm all}
                                                               \tag{E.99.1}
\]
whose abelianized cocycle is zero.  Its origin is the reference
tangential class, and every nonzero member is noncuspidal, including
with respect to the cusp normalizers of this full metabelian quotient.
In particular \(d_3=24\) and \(d_5=168\).
For \(n=3\) there is also a one-parameter family whose class-two
central cohomology coordinate maps to \(9a\kappa_3(3)\) under the
proper exterior-square projection.  For \(a=1\) this central
coordinate is nonzero modulo \(27\).
\end{result}
\begin{proof}
The step retaining mixed-prime finite quotients is an actual module
embedding.  We give its proof.  For a free profinite group \(N\) of
finite rank \(d\), set
\[
 D=N'/N'',\quad R=\widehat{\mathbb Z}[[A]],\quad
 P=N^{(p)},\quad D_p=P'/P'',\quad
 A_p=P^{\rm ab},\quad R_p=\mathbb Z_p[[A_p]].
\]
A free basis gives the continuous Alexander sequences
\[
\begin{aligned}
 0&\longrightarrow D\longrightarrow R^d
          \xrightarrow{\partial}I_A\longrightarrow0,\\
 0&\longrightarrow D_p\longrightarrow R_p^d
          \xrightarrow{\partial_p}I_{A_p}\longrightarrow0,
 \qquad \partial(r_i)=\sum_i r_i(X_i-1).
\end{aligned}                                                   \tag{E.99.2}
\]
Here is a direct justification of the completion in this assertion.
The finite abelian covering graphs of a bouquet of \(d\) circles
have precisely the displayed edge-to-vertex boundary maps with
finite group rings and coefficients.  Their first homology is the
abelianization of the covering subgroup.  These covers detect every
finite character of the closed group \(N'\): given such a character,
choose an open normal \(V_0\trianglelefteq N\) with
\(N'\cap V_0\) in its kernel, and extend the character to \(N'V_0\)
by making it zero on \(V_0\).  This is well-defined since
\([V_0,N']\subseteq N'\cap V_0\).
Conversely, a compatible family of elements in the finite
abelianizations of these covering groups lifts to an element of
\(N'\): the corresponding nonempty closed subsets of compact \(N\)
have the finite intersection property, and their intersection lies
in every preimage of an open subgroup of \(A\), hence in \(N'\).
Taking the inverse limits of the finite graph kernels gives exactly
(E.99.2), not merely its dense image.  The pro-\(p\) proof is the same.

In the coefficient-\(p\) factor of \(R\), average over the
characteristic subgroup \(A^{(p')}=\prod_{q\ne p}A^{(q)}\).
At a finite quotient its normalized sum is an idempotent because its
order is prime to \(p\).  These sums are compatible under pushforward.
Extending the resulting idempotent by zero in the other coefficient
factors gives a central idempotent \(e_p\in R\) with
\[
 e_pR\simeq R_p,\qquad e_pI_A\simeq I_{A_p}.
\]
Comparing kernels in (E.99.2) proves that the natural comparison
\(q:D\to D_p\) restricts to an isomorphism \(e_pD\to D_p\).
Its inverse is therefore a continuous injection
\[
 j_p:D_p\hookrightarrow D,\qquad qj_p=1.                  \tag{E.99.3}
\]
The idempotent is characteristic, and the comparison is natural.
Thus (E.99.3) is equivariant for the full actual arithmetic action.
No group section \(P\to N\) is required or asserted.

Put \(V=P^{\rm ab}\) and \(B=\bigwedge^2_{\mathbb Z_p}V\).
There is an equivariant surjection \(D_p\to B\), the actual
degree-two commutator quotient.  We need surjectivity on continuous
\(H^1\), not only on these modules.  Choose a complement \(\Delta\)
in (E.98.1), using the coprime splitting theorem.  Given a
\(B\)-cocycle on \(\Gamma\), averaging changes it by a coboundary
so that its restriction to \(\Delta\) is zero.  Its restriction to
\(P_{\rm ar}\) is then exactly invariant under
\[
 (a\cdot b)(g)=a\,b(a^{-1}ga),\qquad a\in\Delta.
\]
Lift its values on a free basis of \(P_{\rm ar}\) to \(D_p\).
The universal property of a free pro-\(p\) group, applied to the
pro-\(p\) semidirect product \(D_p\rtimes P_{\rm ar}\), gives a
continuous lifted cocycle.  Its average
\(|\Delta|^{-1}\sum_{a\in\Delta}a\cdot b\) is still a lift and is
exactly invariant.  Extending it by zero on \(\Delta\) proves
\[
 H^1_{\rm cts}(\Gamma,D_p)\twoheadrightarrow
 H^1_{\rm cts}(\Gamma,B).                                \tag{E.99.4}
\]
This works for the infinite compact module \(D_p\) on actual
cocycles; it does not assume surjectivity of a derived inverse limit.

The graph argument in E.98 identifies the right side of (E.99.4)
with \(H^1(G_{\mathbb Q,S},B)\).  At good Frobenius the eigenvalues
on \(V\) have absolute values \(q^{1/2}\) on the proper part and
\(q\) on the boundary part.  Thus the eigenvalues on \(B\) have
absolute values \(q,q^{3/2},q^2\); in particular its invariants
vanish.  Formula (E.98.5) computes its free rank.
There are three real cusps and \((3n-3)/2\) conjugate cusp pairs.
On the boundary module, conjugation is minus the permutation action,
with its diagonal minus line removed.  The proper part contributes
\(g\) eigenvalues of each sign.  Hence
\[
 \dim V^+=\frac{n^2-1}{2},\qquad
 \dim V^-=\frac{n^2+3}{2},\qquad
 \dim B^-=\dim V^+\dim V^-=d_n.                           \tag{E.99.5}
\]
This computes the exact integral cohomology rank without requiring
an integral splitting of the generalized-Jacobian lattice.

Choose cocycles whose classes span a free rank-\(d_n\) submodule of
\(H^1(\Gamma,B)\), and lift them by (E.99.4).  For
\(a\in\mathbb Z_p^{d_n}\), take their indicated linear combination
\(b_a\), inflate to \(G_{\mathbb Q}\), and define
\[
 s_a(g)=j_p(b_a(g))s_O(g).                                 \tag{E.99.6}
\]
Since \(D\) is an abelian normal subgroup and \(j_p\) is equivariant,
this is a continuous homomorphism and a splitting of the full
metabelian extension.  It has literally zero abelian cocycle.
The dependence on \((a,g)\) is jointly continuous.

There is no unidentified geometric conjugacy in this family.
For each coefficient prime \(q\), a good Frobenius different from
\(q\) has no eigenvalue one on \(A^{(q)}\otimes\mathbb Q_q\).
Therefore \(H^0(G_{\mathbb Q},A)=0\).  A conjugator between two
zero-abelian sections must lie in \(D\).  It changes their difference
by a \(D\)-coboundary, whose projection through \(q\) and then to
\(B\) is a coboundary.  Inflation introduces no new coboundaries
because the action already factors through \(\Gamma\).  The
chosen free submodule proves injectivity of (E.99.1).  The compact
parameter space and compact geometric conjugation group also give
the asserted continuity in the section-class topology.

We check the cusp assertion inside \(M_{\rm all}\).  A cusp loop
\(x\) is primitive in a free basis of \(N\), obtained by omitting a
different cusp in a surface presentation.  If an element of
\(M_{\rm all}\) centralizes \(x\), its class-two commutator shows
that its abelian coordinate is a multiple of the primitive vector
\([x]\).  Dividing by that power of \(x\) puts it in \(D\), where
it is annihilated by \(X-1\).  Multiplication by \(X-1\) is injective
on \(R\): in the coefficient-\(q\) factor, refine a finite abelian
quotient by \(q^m\) in the primitive \(X\)-direction.  An invariant
coefficient pushes forward as a multiple of \(q^m\), so vanishes
at every fixed coefficient level in the inverse limit.
The embedding in (E.99.2) gives
\[
 C_{M_{\rm all}}(x)=\overline{\langle x\rangle}.
                                                               \tag{E.99.7}
\]
A geometric normalizer cannot change the exponent of \(x\), since
it fixes its primitive abelian vector.  Its normalizer is therefore
the same group.  Distinct boundary lines in \(A\) distinguish the
cusps, so the arithmetic normalizer of a nonrational cusp cannot
project onto \(G_{\mathbb Q}\).  At a rational cusp the normalizer
has precisely the usual Kummer packet.

For a rational cusp \(d\ne O\), its Abel--Jacobi class is nonzero:
a principal divisor \(d-O\) would give a degree-one map to
\(\mathbb P^1\), contradicting \(g\geq1\).  It is killed by
\(n=p^r\), so its proper \(p\)-adic Kummer class is nonzero by
Mordell--Weil and Kummer theory.  Its packet cannot have zero
abelian class.  At \(O\), the primitive cusp line \(I_O\subset A\)
has \(H^0(G_{\mathbb Q},A/I_O)=0\), again by the proper and boundary
weights.  Hence \(H^1(G_{\mathbb Q},I_O)\to H^1(G_{\mathbb Q},A)\)
is injective.  The only zero-abelian class in that packet is
\([s_O]\).  Injectivity of (E.99.1) now excludes every nonzero
parameter.  This reconstructs LM6 E.33 and E.34 in the full
profinite metabelian quotient.

For the cubic one-parameter assertion we can avoid the unpublished
large matrix certificate of LM6 E.31.  Write \(E=X_3\) with origin
\(O\), and let \(J_D\) be its generalized Jacobian for the nine
cusps.  Its extension class by the boundary torus is killed by
three, because every cusp difference is three-torsion.  Pullback
by \([3]:E\to E\) therefore splits.  The resulting homomorphism
\(j_0:E\to J_D\), with proper projection \([3]\), is unique since
\(\operatorname{Hom}(E,T)=0\), and hence is defined over
\(\mathbb Q\).  On Tate modules it gives an equivariant map
\(T_3E\to V\) whose composite with the proper projection is
multiplication by three.  Its exterior square gives an equivariant
map \(\mathbb Z_3(1)\to B\) whose proper exterior-square projection
is multiplication by nine.  Apply this map to \(\kappa_3(3)\), and
lift the resulting \(B\)-cocycle using (E.99.4).  Formula (E.99.6)
then gives the required one-parameter family.  Its projection is
\(9a\kappa_3(3)\), which has infinite order when \(a\ne0\).
Modulo \(27\), its value as a Kummer class is represented by
\(3^9\), which is not a \(27\)-th power because its valuation at
three is nine.

This verifies the central-cohomology reading of LM6 E.31's
modulo-\(27\) detector.  It does not automatically verify
nonconjugacy in a wholly finite quotient in which \(V\) has also
been reduced modulo \(27\): such a quotient can have nonzero
\(H^0(G_{\mathbb Q},V/27V)\), producing additional conjugators.
The full-metabelian injectivity just proved uses the integral
vanishing \(H^0(G_{\mathbb Q},A)=0\).  No stronger finite-conjugacy
claim is needed here.
\end{proof}

\begin{result}{E.100}{Proved}{The universal Fermat reading is false, and these families are not ordinary tripod pairs}
For every positive even \(n\), the geometrically pro-2 abelian
extension of the plus Fermat curve \(U_n/\mathbb Q\) has no section.
Consequently the same is true for its full geometrically abelian,
metabelian, pronilpotent and profinite extensions.  Thus an
unqualified assertion that every Fermat curve has noncuspidal
sections in all the quotients named in R8 is false.

For \(n=3^r\) or \(5^r\), E.98 and E.99 prove the precise positive
replacement over \(\mathbb Q\).  Every full pronilpotent section
constructed by E.98 has a nonordinary tripod abelian pair at
\(\mathbb Q_2\).  Every full metabelian section in E.99 has tripod
pair \((1,1)\), and is also nonordinary there.  Neither construction
answers the ordinary-pair limit question.
\end{result}
\begin{proof}
If \(n\) is even, \(X_n(\mathbb R)=\varnothing\).  Complex
conjugation acts freely on the punctured complex surface.  Its
quotient is a nonorientable punctured surface with Euler
characteristic \(-n^2/2\), and its fundamental group \(J\) is free
of rank \(1+n^2/2\).  The orientation kernel \(N_{\rm top}\) has
index two.  A free basis of \(J\) can be chosen as
\(t,a_1,\ldots,a_{n^2/2}\), with only \(t\) mapping nontrivially
to \(C_2\); elementary free-basis changes put any nonzero orientation
character in this form.  Schreier's basis for the kernel is
\[
 t^2,\qquad a_i,\qquad ta_it^{-1}.
\]
On its pro-2 abelianization conjugation by \(t\) fixes the first
generator and interchanges each remaining pair.  Thus
\[
 H_1(U_n(\mathbb C),\mathbb Z_2)
   \simeq \mathbb Z_2\oplus
           \mathbb Z_2[C_2]^{n^2/2},\qquad
 H^2(C_2,H_1)\simeq\mathbb Z/2.                           \tag{E.100.1}
\]
The square of the lift \(t\) is the first generator.  Replacing
\(t\) by \(bt\) changes that square by \((1+c)b\), whose first
coordinate is even.  The extension class is therefore the nonzero
class in (E.100.1), and there is no abelian splitting.
The equivariant Riemann-existence comparison identifies the completed
topological extension with the real arithmetic extension; the precise
comparison is recalled in \citet[\S1, equation (4), author p.~3]{LitWickelgrenReal2014}.
We use that comparison, not her theorem requiring real points.
Restriction of a rational section to \(G_{\mathbb R}\) gives the
contradiction.  This independently proves the scope of LM6 E.5.

The recovered review of LM7 E.25 identifies its related quaternion
class as \((-1,-1)\).  Its support includes both two and infinity:
the associated conic \(x^2+y^2+z^2=0\) has no real point, and a
primitive solution over \(\mathbb Z_2\) is impossible modulo four.
At every odd prime the conic has a nonsingular point over the residue
field and lifts by Hensel's lemma.  The phrase ``only at the real
place'' would therefore be false.  None of the positive odd-prime
constructions removes this dyadic obstruction for even \(n\).

We now calculate the positive families' abelian projection.  For
odd \(n\), use the units \(u=-X/Z\), \(v=-Y/Z\).  Then
\[
 t=u^n,\qquad 1-t=v^n,
\]
and (E.98.2) shows that their leading values at the unit tangent
are both one.  The reference tripod pair is consequently \((1,1)\).
The \(p\)-cocycle in E.98 is unramified outside \(p\); its two unit
coordinates lie in
\(H^1(G_{\mathbb Q,S},\mathbb Z_p(1))=\mathbb Z_p\kappa_p(p)\).
Write them as \(a\kappa_p(p),b\kappa_p(p)\).  If
\(\epsilon_p\in\widehat{\mathbb Z}=\prod_q\mathbb Z_q\) is the
idempotent supported at \(p\), the actual completed tripod pair is
\[
             \bigl(p^{na\epsilon_p},p^{nb\epsilon_p}\bigr).
                                                               \tag{E.100.2}
\]
This equality includes every coefficient prime.  In particular its
2-primary component is \((1,1)\).  The homomorphism
\[
 \mathbb Q_2^\times\longrightarrow
       \varprojlim_m\mathbb Q_2^\times/\mathbb Q_2^{\times 2^m}
                                                               \tag{E.100.3}
\]
is injective, as is immediate from
\(\mathbb Q_2^\times=2^{\mathbb Z}\times\{\pm1\}\times
(1+4\mathbb Z_2)\).
An ordinary local point with pair (E.100.2) would therefore satisfy
\(t=1\) and \(1-t=1\) in \(\mathbb Q_2^\times\), a contradiction.
The full metabelian family has zero intrinsic abelian cocycle, so its
image in the rank-two tripod abelianization is exactly the reference
pair \((1,1)\), giving the same contradiction.

Thus the word ``noncuspidal'' in the recovered positive Fermat
theorems cannot be replaced by ``locally ordinary on the tripod''.
The individual projections of the E.98 product section at
\(q\ne p\) are deliberately cuspidal.  Noncuspidality of the full
pronilpotent section consequently does not assert noncuspidality of
each of its coefficient-prime projections.
The full pronilpotent quotient in E.98 and the full metabelian
quotient in E.99 are also different quotient constructions; the
Alexander-module embedding is the additional argument making the
latter retain mixed-prime finite quotients.  Finally, Hoshi's
published Theorem B has its own field and regular-prime hypotheses,
including \(p>7\) \citep[author p.~5]{LitHoshi2010}.  It is not
being cited as a statement of the above \(\mathbb Q\), \(p=3,5\)
families.  Their derivation has been given here.
\end{proof}

\paragraph{Recovered and missing historical pieces.}
The full original LM6 proof file, including E.16, E.31--E.34 and
E.38, was recovered.  Its referenced finite scripts
\texttt{full\_metabelian\_family.py} and
\texttt{metabelian\_family\_rank.py} were not present; E.99 gives
independent replacements for the conclusions used here.  The original
LM7 files named in the recovered review---in particular
\texttt{E\_completed\_pairs.tex}, \texttt{E\_cubic.tex} and
\texttt{E\_intrinsic\_positive.tex}---and the original LM8 F.21
proof have not been recovered.  The \texttt{LM7\_PUSH} checkpoint
used in E.57 is a different run and does not repair that source gap.
A.24 independently proves the derived-limit phenomenon in R8; it
does not certify every marked geometric coordinate of LM6 E.38.
That latter statement additionally uses E.19's precise obstruction,
E.37's integral elliptic action map, and a missing 330-dimensional
Smith certificate.  Those extra geometric and numerical assertions
are not used in E.98--E.100 and remain separately unaudited.
